\section{Oracle experiment plan}
The oracle study removes learning imperfections deliberately. True family labels
are used, and long low-noise datasets support high-quality individual LPV models.
The comparison in \cref{tab:methods} isolates scheduling and specialization.

\begin{table}[ht]
\centering
\caption{Core oracle model architectures.}
\label{tab:methods}
\begin{tabular}{llllr}
\toprule
Method & Scheduling & Specialization & Interpretation & Families\\
\midrule
Global LTI & none & none & ignores both factors & 1\\
Clustered LTI & none & oracle family & structural effect only & 3\\
Global LPV & continuous & none & scheduling effect only & 1\\
Clustered LPV & continuous & oracle family & both effects & 3\\
Gridded LTI & discrete & family and speed & high-complexity oracle & 15\\
\bottomrule
\end{tabular}
\end{table}

\subsection{Execution phases}
\begin{enumerate}
\item \textbf{Plant validation.} Inspect continuous and frozen discrete dynamics,
  stability, handling properties, family separation, and within-family scatter.
\item \textbf{LPV approximation validation.} Fit the proposed basis and measure
  matrix approximation errors on a fine speed grid. Revise the basis if needed.
\item \textbf{Oracle model comparison.} Evaluate one-step prediction, matrix error,
  and generalization for the four factorial methods and the gridded oracle.
\item \textbf{Controller comparison.} Use common LQI weights and gain scheduling on
  fixed-family/changing-speed, fixed-speed/different-family, combined, and fully
  unseen test scenarios.
\item \textbf{LTI complexity sweep.} Sweep $K_{\mathrm{LTI}}\in\{1,2,3,5,8,10,15\}$
  and determine how many prototypes are needed to match three LPV families.
\item \textbf{Two-dimensional regime study.} Sweep structural heterogeneity
  $\gamma_h\in\{0,0.5,1,1.5\}$ and narrow, medium, and broad speed envelopes.
\end{enumerate}

\subsection{Metrics and scenarios}
Model quality is assessed through one-step state prediction, separate $\beta$ and
$r$ RMSE, and $\lVert A_i(\rho)-\hat A(\rho)\rVert_F$ and
$\lVert B_i(\rho)-\hat B(\rho)\rVert_F$ versus speed. Closed-loop evaluation uses
yaw-rate tracking RMSE, sideslip RMS, steering RMS and peaks, steering-rate peaks,
and a frozen-grid spectral-radius audit. Every result reports fleet average,
dispersion, worst-family performance, and the family performance gap. Ten seeds
vary client realizations and test trajectories.

\subsection{Predeclared decision gates}
The idea passes the oracle gate only if: (i) global LPV materially improves over
global LTI over the broad speed envelope; (ii) clustered LPV improves over global
LPV, particularly for the worst family; (iii) clustered LPV improves over
clustered LTI on unseen varying-speed tests; (iv) three LPV families match or beat
substantially more than three LTI prototypes; and (v) benefits generalize to
unseen speed--maneuver combinations. Working thresholds are more than 10\%
improvement for LPV relevance and approximately 5--10\% for specialization, but
effect consistency and practical scale take precedence over arbitrary cutoffs.

The project is paused or reformulated if global and clustered LPV are equivalent,
clustered LTI matches clustered LPV, a few LTI prototypes match the LPV solution,
benefits disappear out of distribution, or only unrealistic family separation
creates the result.

\subsection{Planned evidence package}
The minimum evidence consists of six figures: ground-truth model entries versus
speed, the LPV manifolds in model space, model error versus speed, a representative
closed-loop trajectory, performance versus model/controller complexity, and a
two-dimensional operating-range--heterogeneity regime map.

